\documentclass[11pt]{article}

\usepackage[margin=2.5cm]{geometry}
\usepackage{amsmath,amssymb,amsfonts,amsthm}
\usepackage{graphicx}
\usepackage{booktabs}
\usepackage{algorithm}
\usepackage{algpseudocode}
\usepackage{hyperref}
\usepackage[numbers,sort&compress]{natbib}
\usepackage{xcolor}
\usepackage{microtype}
\usepackage{caption}
\usepackage{listings}
\usepackage{xspace}

% Custom commands (mirrors main.tex)
\newcommand{\falcon}{\textsc{Falcon}\xspace}
\newcommand{\fastprop}{\textsc{FastProp}\xspace}
\newcommand{\randprop}{\textsc{RandProp}\xspace}
\newcommand{\crossnet}{\textsc{CrossNet}\xspace}
\newcommand{\Cov}{\mathrm{Cov}}
\newcommand{\Var}{\mathrm{Var}}
\newcommand{\clr}{\mathrm{clr}}
\newcommand{\R}{\mathbb{R}}
\newcommand{\bfX}{\mathbf{X}}
\newcommand{\bfY}{\mathbf{Y}}
\renewcommand{\thesection}{S\arabic{section}}
\renewcommand{\thefigure}{S\arabic{figure}}
\renewcommand{\thetable}{S\arabic{table}}
\renewcommand{\theequation}{S\arabic{equation}}

\lstset{
  basicstyle=\ttfamily\footnotesize,
  breaklines=true,
  frame=single,
  rulecolor=\color{gray!50},
  numbers=left,
  numbersep=6pt,
  numberstyle=\tiny\color{gray},
  keywordstyle=\color{blue!70!black},
  language=Python,
}

\theoremstyle{plain}
\newtheorem{proposition}{Proposition}
\newtheorem{lemma}{Lemma}

\title{Supplementary Information for\\
\falcon: Fast Algorithm for Large-scale Cross-domain\\
Compositional Network Inference}
\author{}
\date{}

\begin{document}
\maketitle

\tableofcontents
\vspace{1em}

% ============================================================
\section{Derivation of the cross-domain centering identity}
\label{sec:derivation}

For independently normalised compositions $\bfX \in \Delta^{p-1}$ and
$\bfY \in \Delta^{q-1}$, let $W_i^X, W_j^Y$ denote the underlying absolute
abundances with totals $S^X = \sum_k W_k^X$ and $S^Y = \sum_l W_l^Y$. After
the centered log-ratio transform,
\begin{equation}
\clr_X(i) \;=\; \log W_i^X - \tfrac{1}{p}\sum_{k=1}^{p}\log W_k^X
\quad\text{(and analogously for $Y$).}
\end{equation}
Define the true cross-covariance of the log-abundances,
$\Omega_{ij} \;\triangleq\; \Cov(\log W_i^X, \log W_j^Y)$.
Substituting and using bilinearity of covariance,
\begin{align}
T_{ij} &\triangleq \Cov(\clr_X(i),\, \clr_Y(j)) \nonumber \\
       &= \Omega_{ij}
        - \tfrac{1}{p}\sum_{k} \Omega_{kj}
        - \tfrac{1}{q}\sum_{l} \Omega_{il}
        + \tfrac{1}{pq}\sum_{k,l}\Omega_{kl}. \label{eq:centering-scalar}
\end{align}
The four terms have a compact matrix interpretation: each is a different
component of a double-centered $\Omega$.

\begin{proposition}[Matrix form of the centering bias]\label{prop:matrix}
Let $H_p = I_p - \tfrac{1}{p}\mathbf{1}_p\mathbf{1}_p^\top$ denote the
$p$-dimensional centering matrix. Then
$T = H_p\,\Omega\,H_q^\top$.
\end{proposition}

\begin{proof}
$(H_p \Omega H_q^\top)_{ij}
= \Omega_{ij} - \tfrac{1}{p}\sum_k\Omega_{kj} - \tfrac{1}{q}\sum_l\Omega_{il}
+ \tfrac{1}{pq}\sum_{k,l}\Omega_{kl}$,
which matches \eqref{eq:centering-scalar} term-by-term. The two outer sums
expand to the marginal and double-marginal corrections; the inner product
gives the bilinear term.
\end{proof}

\begin{lemma}[Non-identifiability without regularisation]\label{lem:nonid}
For any constants $a \in \R^p$ and $b \in \R^q$,
$T(\Omega + \mathbf{1}_p a^\top + b\,\mathbf{1}_q^\top) = T(\Omega)$.
\end{lemma}

\begin{proof}
$H_p \mathbf{1}_p = 0$ and $H_q^\top \mathbf{1}_q = 0$ because the all-ones
vector lies in the null space of any centering matrix. Hence
$H_p(\mathbf{1}_p a^\top)H_q^\top = 0$, and the same for the $b$ shift.
\end{proof}

Lemma~\ref{lem:nonid} formalises the central identifiability obstacle: the
$(p+q-1)$-dimensional null space of $(H_p, H_q^\top)$ leaves $\Omega$
underdetermined. The $\ell_1$ penalty in equation~(7) of the main text breaks
this degeneracy by selecting the unique solution of minimum mass among the
infinite family compatible with $T$, encoding the biological prior that most
phage--bacteria pairs are non-interacting.

% ============================================================
\section{FISTA convergence and step size}
\label{sec:fista}

Let $f(\Omega) = \tfrac{1}{2}\|T - H_p\Omega H_q^\top\|_F^2
+ \tfrac{\lambda_2}{2}\|\Omega - \Omega_{\mathrm{prior}}\|_F^2$. Using the
idempotence $H_p^2 = H_p$ and symmetry $H_p = H_p^\top$,
\begin{equation}
\nabla f(\Omega) = H_p\,(H_p \Omega H_q^\top - T)\,H_q
                 + \lambda_2 (\Omega - \Omega_{\mathrm{prior}}).
\end{equation}

\begin{proposition}[Lipschitz constant]
$\nabla f$ is Lipschitz with constant $L \leq 1 + \lambda_2$.
\end{proposition}

\begin{proof}
The spectral norm of any centering matrix satisfies $\|H_p\|_2 = 1$ because
its eigenvalues lie in $\{0, 1\}$. Hence the linear map
$\Omega \mapsto H_p \Omega H_q^\top$ has operator norm bounded by
$\|H_p\|_2 \|H_q\|_2 = 1$, and the gradient of the Frobenius residual is
$1$-Lipschitz. Adding the quadratic Tikhonov term contributes $\lambda_2$.
\end{proof}

We use step size $\eta = 1/L = 1/(1 + \lambda_2)$. With Nesterov momentum
\citep{Nesterov1983,BeckTeboulle2009}, FISTA achieves
$F(\Omega_k) - F(\Omega^\star) = O(1/k^2)$, where
$F(\Omega) = f(\Omega) + \lambda_1\|\Omega\|_1$. In the simulations of
Section~3.5 of the main text, the relative change
$\|\Omega_{k+1} - \Omega_k\|_F / \|\Omega_k\|_F$ drops below $10^{-4}$
within $20$--$30$ iterations.

% ============================================================
\section{Adaptive regularisation parameters}
\label{sec:lambda}

The sparsity parameter $\lambda_1$ controls the trade-off between recall
and precision. We default to the BIC-style choice
\begin{equation}
\lambda_1 \;=\; \kappa \cdot \hat{\sigma}_T \sqrt{\log(pq)/n},
\end{equation}
where $\hat{\sigma}_T$ is the standard deviation of the entries of the
observed $T$ matrix and $\kappa$ is a sparsity prior in $(0,1)$ controlling
how aggressively to threshold. The scaling $\sqrt{\log(pq)/n}$ matches the
near-optimal rate for sparse signal detection in regularised regression.
Practitioners with stability-based criteria (e.g.\ StARS \citep{Liu2010})
may replace this default by repeatedly subsampling the data and choosing the
largest $\lambda_1$ at which the empirical edge variability is below a
nominal $\beta = 0.05$.

The prior weight $\lambda_2$ is typically small ($10^{-2}$ to $10^{-1}$
times $\|T\|_F^2/(pq)$) when biological priors are partial or uncertain.
Setting $\lambda_2 = 0$ disables the prior entirely.

% ============================================================
\section{Multiplicative replacement and the CLR bias}
\label{sec:zero-handling}

The CLR transform requires strictly positive entries. We follow
\cite{MartinFernandez2003}:
\begin{equation}
x_i^{\mathrm{rep}} =
\begin{cases}
\delta & \text{if } x_i = 0, \\[2pt]
x_i \cdot \left(1 - n_0 \delta\right) & \text{otherwise,}
\end{cases}
\qquad
\delta = \frac{0.65}{p^2},
\end{equation}
where $n_0$ is the number of zero entries in the row. This preserves the
unit sum constraint exactly. The residual bias introduced by zero
replacement scales as $O(1/p^2)$ in the CLR coordinates (cf.\ Lemma~3 of
\cite{Quinn2019}) and is therefore negligible in the regimes we target
($p \geq 50$).

% ============================================================
\section{Random projection: dimension selection}
\label{sec:randprop}

\randprop uses the Johnson--Lindenstrauss lemma
\citep{JohnsonLindenstrauss1984,Achlioptas2003} to embed the
column-normalised CLR matrix $\tilde Z \in \R^{n \times p}$ into a
$d$-dimensional sketch. For pairwise inner-product preservation within
tolerance $\varepsilon$ with probability $\geq 1 - 1/p^2$, we set
\begin{equation}
d = \left\lceil 8 \ln(p) / \varepsilon^2 \right\rceil,
\end{equation}
clipped to $\min(d, n)$ since exceeding the sample dimension yields no
additional information. The Achlioptas sparse construction draws each entry
of the projection matrix from $\{+1/\sqrt{d/3}, 0, -1/\sqrt{d/3}\}$ with
probabilities $\{1/6, 2/3, 1/6\}$, reducing the effective multiplication
count by a factor of three without sacrificing the JL guarantee.

The top-$k$ extraction step uses \texttt{numpy.argpartition} (introselect)
to obtain the $k$ largest absolute correlations per row in $O(p)$ per row.
Refinement recomputes the exact $\rho_p$ for the candidate set in
$O(nk)$ per row.

% ============================================================
\section{Extended simulation tables}
\label{sec:tables}

\begin{table}[h]
\centering
\caption{Single-domain detection power at $p = 100$
(mean $\pm$ standard deviation across 20 replicates). Power is the fraction
of true edges detected at BH-FDR $< 0.05$ with a minimum effect-size filter
of $|\rho_p| > 0.1$.}
\label{tab:power-p100}
\begin{tabular}{@{}cccccc@{}}
\toprule
$\rho$ \textbackslash\ $n$ & 50 & 100 & 200 & 500 \\
\midrule
0.2 & $0.00\pm0.00$ & $0.00\pm0.00$ & $0.00\pm0.00$ & $0.00\pm0.00$ \\
0.3 & $0.00\pm0.00$ & $0.00\pm0.00$ & $0.00\pm0.00$ & $0.58\pm0.09$ \\
0.4 & $0.00\pm0.00$ & $0.00\pm0.00$ & $0.00\pm0.00$ & $0.99\pm0.01$ \\
0.5 & $0.00\pm0.00$ & $0.00\pm0.00$ & $0.01\pm0.02$ & $1.00\pm0.00$ \\
0.7 & $0.00\pm0.00$ & $0.00\pm0.00$ & $0.66\pm0.05$ & $1.00\pm0.00$ \\
\bottomrule
\end{tabular}
\end{table}

\begin{table}[h]
\centering
\caption{Single-domain detection power at $p = 500$. The conservative
behaviour reflects BH-FDR adjustment over
$\binom{500}{2} \approx 1.25\times 10^{5}$ tests, not an intrinsic failure
of the proportionality estimator: pairwise $|\rho_p|$ estimates remain
accurate (Pearson correlation with truth $> 0.95$ in all conditions; not
shown). \randprop's top-$k$ pre-selection or stability-based
thresholding \citep{Liu2010} is required at this scale.}
\label{tab:power-p500}
\begin{tabular}{@{}cccccc@{}}
\toprule
$\rho$ \textbackslash\ $n$ & 50 & 100 & 200 & 500 \\
\midrule
0.2 & 0.00 & 0.00 & 0.00 & 0.00 \\
0.3 & 0.00 & 0.00 & 0.00 & 0.00 \\
0.4 & 0.00 & 0.00 & 0.00 & 0.00 \\
0.5 & 0.00 & 0.00 & 0.00 & 0.00 \\
0.7 & 0.00 & 0.00 & 0.00 & 0.00 \\
\bottomrule
\end{tabular}
\end{table}

\begin{table}[h]
\centering
\caption{Cross-domain inference at $p = 60$, $q = 80$, $n = 300$,
20 planted interactions (70\% lytic / negative), 20 replicates.
Significance threshold $|C| > 0.1$ and Fisher-$z$ $p < 0.05$.}
\label{tab:cross-domain}
\begin{tabular}{@{}lccccc@{}}
\toprule
Method & Correlation & Sign acc. & Sensitivity & Specificity & Bias \\
\midrule
Naive CLR & $0.993\pm 0.005$ & $1.000$ & $1.000$ & $0.9489\pm 0.004$ & $\hphantom{-}\,0.0006$ \\
SparXCC-like & $0.993\pm 0.005$ & $1.000$ & $1.000$ & $0.9489\pm 0.004$ & $-0.00005$ \\
\crossnet & $0.993\pm 0.005$ & $1.000$ & $1.000$ & $\mathbf{0.9622\pm 0.003}$ & $-0.00009$ \\
\bottomrule
\end{tabular}
\end{table}

\begin{table}[h]
\centering
\caption{Wall-clock benchmarks on a single CPU core (Apple Silicon,
NumPy 1.26 with Accelerate BLAS), median of 3 runs.}
\label{tab:scalability}
\begin{tabular}{@{}rrrr@{}}
\toprule
$n$ & $p$ & \fastprop\ (s) & \randprop\ (s) \\
\midrule
100 & 100 & 0.0002 & 0.007 \\
100 & 500 & 0.002 & 0.022 \\
100 & 1000 & 0.008 & 0.045 \\
100 & 2000 & 0.040 & 0.157 \\
100 & 5000 & 0.315 & 0.366 \\
500 & 500 & 0.003 & 0.061 \\
500 & 1000 & 0.010 & 0.117 \\
500 & 2000 & 0.043 & 0.259 \\
500 & 5000 & 0.420 & 1.696 \\
1000 & 2000 & 0.052 & 0.697 \\
\bottomrule
\end{tabular}
\end{table}

% ============================================================
\section{Reproducibility checklist}
\label{sec:repro}

\begin{itemize}
\item Algorithms: \texttt{src/falcon/\_\_init\_\_.py} (Python 3.10+,
  NumPy~$\geq$~1.24, SciPy~$\geq$~1.10).
\item Benchmark runner: \texttt{benchmarks/run\_on\_server.py} with
  \texttt{multiprocessing.Pool}, regenerating every CSV table in
  \texttt{data/} (\texttt{scalability.csv}, \texttt{detection.csv},
  \texttt{cross\_domain.csv}, \texttt{fdr\_control.csv}) from documented
  random seeds.
\item CSV I/O helpers: \texttt{benchmarks/io\_utils.py}.
\item Figure generation: \texttt{manuscript/figures/generate\_workflow.py},
  \texttt{generate\_single\_domain.py}, \texttt{generate\_cross\_domain.py}.
  Each script reads only the CSV tables in \texttt{data/} so figures can
  be regenerated without re-running the benchmarks.
\item Repository will be archived on Zenodo at the time of acceptance with
  a permanent DOI; see Data and Code Availability.
\end{itemize}

\bibliographystyle{unsrtnat}
\bibliography{../references}

\end{document}
